% element: 10-s004:e04
\BeleskeID{10-s004}
\subsection*{Od kašnjenja do provođenja}
U intervalu $t_0$–$t_1$ puni se oblast prostornog tovara, odnosno njena ekvivalentna kapacitivnost. Oblast koja je bila inverzno polarizovana skuplja se, ali spoj je još uvek inverzno polarizovan. To je vreme kašnjenja.

Pri kraju tog intervala pn spoj postaje direktno polarizovan. U $t_1$–$t_2$ dioda vodi, a struja je određena veličinom ulaznog napona, otpornošću i direktnim naponom diode:
\[I_D=\frac{V_H-V_D}{R}.\]
\subsection*{Isključivanje i oporavak}
U $t_2$–$t_3$ postoje nagomilani nosioci uz oblast prostornog tovara. Spoj je i dalje direktno polarizovan iako je ulaz promenjen. Negativna struja odnosi te nosioce:
\[I_D=\frac{V_L-V_D}{R},\qquad V_L<0\ \Longrightarrow\ I_D<0.\]
U intervalu $t_3$–$t_4$ pn spoj postaje inverzno polarizovan i oblast prostornog tovara se širi. Promena naelektrisanja ekvivalentnih kapacitivnosti daje eksponencijalan oblik prelaza.
